\appendix

\section*{Appendix A: Review of Differential Geometry}
\addcontentsline{toc}{section}{Appendix A: Review of Differential Geometry}

\section{Manifolds and Smooth Maps}

A \emph{smooth manifold} of dimension $n$ is a second countable Hausdorff space $M$ together with an atlas of coordinate charts ${(U_i,\varphi_i)}$ where $\varphi_i\colon U_i\to\mathbb{R}^n$ are homeomorphisms onto their images and all transition maps $\varphi_j\circ\varphi_i^{-1}$ are smooth.

\begin{definition}
A map $f\colon M\to N$ between smooth manifolds is called smooth if for every pair of charts $\varphi$ on $M$ and $\psi$ on $N$, the coordinate representation $\psi\circ f\circ\varphi^{-1}$ is a smooth function between Euclidean spaces.
\end{definition}

Note that in the categorical viewpoint, manifolds form a subcategory of smooth schemes, and ultimately of derived stacks.  The lamphrodynamic formalism requires a broad generalization of these notions, but ordinary smooth manifolds remain the geometric base of physical spacetime throughout Volume~I.

\subsection{Tangent Bundle}

For each $p\in M$, the tangent space $T_pM$ is defined as equivalence classes of smooth curves passing through $p$, or equivalently as derivations on germs of smooth functions at $p$.  The union of all tangent spaces
[
TM=\bigsqcup_{p\in M}T_pM
]
is a smooth vector bundle over $M$.

\begin{remark}
Later, tangent complexes of derived stacks generalize these tangent spaces.  In the derived setting, the tangent object is a chain complex rather than an ordinary vector space.
\end{remark}

\section{Vector Bundles and Differential Forms}

A (real) vector bundle over $M$ is a smooth manifold $E$ together with a smooth projection $\pi\colon E\to M$ such that for each $p\in M$, the fiber $E_p=\pi^{-1}(p)$ is a real vector space and locally $E\simeq U\times\mathbb{R}^k$.

\begin{definition}
The space of smooth sections of $E$ is denoted $\Gamma(M,E)$.
\end{definition}

In particular, the cotangent bundle $T^*M$ yields the spaces of differential $k$-forms
[
\Omega^k(M)=\Gamma(M,\wedge^kT^*M).
]

\section{Connections and Curvature}

Let $E\to M$ be a vector bundle. A connection is a $\mathbb{R}$--linear operator
[
\nabla\colon \Gamma(M,E)\to\Gamma(M,T^*M\otimes E)
]
satisfying the Leibniz rule $\nabla(fs)=df\otimes s + f,\nabla s$ for $f\in C^\infty(M)$ and $s\in\Gamma(M,E)$.

\begin{definition}
The curvature of a connection $\nabla$ is the bundle map
[
R_\nabla\colon \Gamma(M,E)\to \Gamma(M,\wedge^2T^*M\otimes E)
]
defined by $R_\nabla(s)=\nabla^2s$.
\end{definition}

Connections appear in RSVP as part of geometric background data, but the lamphrodynamic fields themselves are defined intrinsically from the derived moduli stack, not imposed externally.

\section{Riemannian and Lorentzian Geometry}

A Riemannian metric $g$ on $M$ is a smooth section of $T^*M\otimes T^*M$ which is symmetric and positive definite.  A Lorentzian metric has signature $(-,+,\ldots,+)$ and is the standard setting for general relativity.

The Levi--Civita connection is the unique torsion--free connection compatible with a given metric.  Its curvature tensor and associated invariants (Ricci curvature, scalar curvature) appear throughout classical field theory and reappear in the comparison sections of Chapter~14.

\begin{remark}
One conceptual theme of RSVP is that the lamphrodynamic metric need not be an externally fixed Lorentzian structure. Instead, the triple $(\Phi,\mathbf v,S)$ plays the role of the fundamental geometric data.
\end{remark}

\section{Exterior Calculus}

The exterior derivative $d:\Omega^k(M)\to\Omega^{k+1}(M)$ satisfies $d^2=0$, and the wedge product makes $\Omega^\bullet(M)$ into a graded commutative algebra.  Many constructions in the $(-1)$--shifted symplectic setting rely on derived analogues of differential forms.

\begin{definition}
The de,Rham cohomology of $M$ is the cohomology of the complex $(\Omega^\bullet(M),d)$.
\end{definition}

\begin{remark}
For derived stacks, the de,Rham complex is upgraded to a mixed complex whose cohomology yields the shifted symplectic form used in the AKSZ construction.
\end{remark}

\section{End of Appendix A}

This appendix only summarizes the necessary smooth differential geometry. The derived generalizations in Chapters~2--4 systematically replace tangent spaces by tangent complexes, replace manifolds by derived Artin stacks, and replace ordinary forms by their graded, shifted, or mixed counterparts.

